The \code{BPatch} class represents the entire Dyninst library. There can only be one instance
of this class at a time. This class is used to perform functions and obtain information that
is not specific to a particular thread or image.

\begin{apient}
std::vector<BPatch_process*> *getProcesses()
\end{apient}

\apidesc{
Returns the list of processes that are currently defined. This list includes processes
that were directly created by calling processCreate/processAttach, and
indirectly by the UNIX fork or the Windows CreateProcess system call. It is up to
the user to delete this vector when they are done with it.
}

\begin{apient}
BPatch_process *processAttach(const char *path, int pid, BPatch_hybridMode mode=BPatch_normalMode)

BPatch_process *processCreate(const char *path, const char *argv[], const char **envp = NULL,
                              int stdin_fd=0, int stdout_fd=1, int stderr_fd=2,
                              BPatch_hybridMode mode=BPatch_normalMode)
\end{apient}

\apidesc{
Each of these functions returns a pointer to a new instance of the
\code{BPatch\_process} class. The path parameter needed by these functions should be the
pathname of the executable file containing the process image. The \code{processAttach}
function returns a \code{BPatch\_process} associated with an existing process. On Linux platforms
the path parameter can be NULL since the executable image can be derived from the
process pid. Attaching to a process puts it into the stopped state. The processCreate
function creates a new process and returns a new \code{BPatch\_process} associated with it.
The new process is put into a stopped state before executing any code.

The \code{stdin\_fd}, \code{stdout\_fd}, and \code{stderr\_fd} parameters are used to set the standard
input, output, and error of the child process. The default values of these parameters
leave the input, output, and error to be the same as the mutator process. To
change these values, an open UNIX file descriptor (see \code{open(1)}) can be passed.

The mode parameter is used to select the desired level of code analysis. Activating
hybrid code analysis causes Dyninst to augment its static analysis of the code with
run-time code discovery techniques. There are three modes: \code{BPatch\_normalMode},
\code{BPatch\_exploratoryMode}, and \code{BPatch\_defensiveMode}. Normal mode enables the
regular static analysis features of Dyninst. Exploratory mode and defensive mode
enable addtional dynamic features to correctly analyze programs that contain uncommon
code patterns, such as malware. Exploratory mode is primarily oriented towards analyzing
dynamic control transfers, while defensive mode additionally aims to tackle code
obfuscation and self-modifying code. Both of these modes are still experimental and
should be used with caution. Defensive mode is only supported on Windows. Defensive
mode has been tested on normal binaries (binaries that run correctly under
normal mode), as well as some simple, packed executables (self-decrypting or
decompressing). More advanced forms of code obfuscation, such as self-modifying code,
have not been tested recently. The traditional Dyninst interface may be used for
instrumentation of binaries in defensive mode, but in the case of highly
obfuscated code, this interface may prove to be ineffective due to the lack of a
complete view of control flow at any given point. Therefore, defensive mode also
includes a set of callbacks that enables instrumentation to be performed as new code
is discovered. Due to the fact that recent efforts have focused on simpler forms
of obfuscation, these callbacks have not been tested in detail. The next release
of Dyninst will target more advanced uses of defensive mode.
}

\begin{apient}
BPatch_binaryEdit *openBinary(const char *path, bool openDependencies = false)

\end{apient}

\apidesc{
This function opens the executable file or library file pointed to by path for
binary rewriting. If openDependencies is true then Dyninst will also open all shared
libraries that path depends on. Upon success, this function returns a new instance
of a \code{BPatch\_binaryEdit} class that represents the opened file and any dependent shared
libraries. This function returns NULL in the event of an error.
}

\begin{apient}
bool pollForStatusChange()

\end{apient}

\apidesc{
This is useful for a mutator that needs to periodically check on the status of its managed threads and does not want to
check each process individually. It returns true if there has been a change in the status of one or more threads that
has not yet been reported by either isStopped or isTerminated.
}

\begin{apient}
void setDebugParsing (bool state)

\end{apient}

\apidesc{
Turn on or off the parsing of debugger information. By default, the debugger information (produced by the –g compiler
option) is parsed on those platforms that support it. However, for some applications this information can be quite
large. To disable parsing this information, call this method with a value of false prior to creating a process.
}

\begin{apient}
bool parseDebugInfo()

\end{apient}

\apidesc{
Return true if debugger information parsing is enabled, or false otherwise.
}

\begin{apient}
void setTrampRecursive (bool state)

\end{apient}

\apidesc{
Turn on or off trampoline recursion. By default, any snippets invoked while another snippet is active will not be
executed. This is the safest behavior, since recursively-calling snippets can cause a program to take up all
available system resources and die. For example, adding instrumentation code to the start of printf, and then calling
printf from that snippet will result in infinite recursion.
This protection operates at the granularity of an instrumentation point. When snippets are first inserted at a point,
this flag determines whether code will be created with recursion protection. Changing the flag is not retroactive,
and inserting more snippets will not change the recursion protection of the point. Recursion protection increases the
overhead of instrumentation points, so if there is no way for the snippets to call themselves, calling this method with
the parameter true will result in a performance gain. The default value of this flag is false.
}

\begin{apient}
bool isTrampRecursive ()

\end{apient}

\apidesc{
Return whether trampoline recursion is enabled or not. True means that it is enabled.
}

\begin{apient}
void setTypeChecking(bool state)

\end{apient}

\apidesc{
Turn on or off type-checking of snippets. By default type-checking is turned on, and an attempt to create a snippet
that contains type conflicts will fail. Any snippet expressions created with type-checking off have the type of their
left operand. Turning type-checking off, creating a snippet, and then turning type-checking back on is similar to the
type cast operation in the C programming language.
}

\begin{apient}
bool isTypeChecked()

\end{apient}

\apidesc{
Return true if type-checking of snippets is enabled, or false otherwise.
}

\begin{apient}
bool waitForStatusChange()

\end{apient}

\apidesc{
This function waits until there is a status change to some thread that has not yet been reported by either isStopped or
isTerminated, and then returns true. It is more efficient to call this function than to call pollForStatusChange in a
loop, because waitForStatusChange blocks the mutator process while waiting.
}


\begin{apient}
void addAnalysisExcludePattern(const char *pattern)

\end{apient}

\apidesc{
Exclude shared objects from analysis. A shared object whose name matches \code{pattern} is
still loaded --- its symbol table is read, and its symbols, address ranges and inferior heaps
remain available --- but Dyninst builds no control flow graph for it, so it contains no
functions and no basic blocks. This avoids the cost of parsing libraries that will never be
instrumented, which can dominate process creation for programs that load very large shared
libraries.

\code{pattern} is a shell-style wildcard expression as accepted by \code{fnmatch(3)},
supporting \code{*}, \code{?} and \code{[...]}. It is matched against both the object's full
path and its base name, so \code{libLLVM.so*} matches
\code{/usr/lib64/libLLVM.so.23.0git}. This function may be called more than once; an object is
excluded if it matches any of the patterns.

Patterns must be added before the process is created or the binary is opened, since an object
is analyzed at the point it is loaded. The executable and the Dyninst runtime library are
never excluded.

Dyninst itself resolves a handful of functions by name during process creation, among them
\code{pthread\_self} in \code{libpthread}. Excluding a library that Dyninst needs this way can
make process creation fail; it reports the failure rather than continuing silently.

Excluding a library the mutator later queries is not an error, but lookups into it find
nothing: its functions cannot be found by name or by address, \code{BPatch\_object::findFunction}
returns no match, and per-module queries such as \code{BPatch\_module::getProcedures} come back
empty. Instrumenting an excluded library is therefore not possible.
}

\begin{apient}
void clearAnalysisExcludePatterns()

\end{apient}

\apidesc{
Discard every pattern previously passed to \code{addAnalysisExcludePattern}. Objects that have
already been loaded are unaffected.
}

\begin{apient}
bool analysisExcluded(const char *name)

\end{apient}

\apidesc{
Return true if \code{name} matches any pattern passed to \code{addAnalysisExcludePattern},
using the same matching rules.
}

\begin{apient}
typedef bool (*BPatchAnalyzeObjectCallback)(Dyninst::SymtabAPI::Symtab *symtab)

BPatchAnalyzeObjectCallback
registerAnalyzeObjectCallback(BPatchAnalyzeObjectCallback func)
\end{apient}

\apidesc{
Register a callback that decides whether a shared object is analyzed, and return the
previously registered callback, or NULL. Where \code{addAnalysisExcludePattern} matches on a
name known in advance, this decides from the object's contents: the callback is passed that
object's \code{Symtab}, so it can consult the size of the symbol table --
\code{getAllFunctionsRef().size()} -- or anything else SymtabAPI exposes. Skipping every
library above some number of functions, for instance, cannot be written as a pattern.

The callback is invoked once per shared object as the object is loaded, after its symbol table
has been read but before any control flow graph is built. Returning false leaves the object
loaded but unparsed, exactly as a matching pattern would; returning true analyzes it normally.

It is not consulted for the executable, for the Dyninst runtime library, or for an object that
a pattern has already excluded. It must be registered before the process is created or the
binary is opened, and, because it runs inside the load path, it must not call back into
Dyninst.
}

\begin{apient}
BPatchAnalyzeObjectCallback getAnalyzeObjectCallback()
\end{apient}

\apidesc{
Return the currently registered analysis callback, or NULL if none has been registered.
}

\begin{apient}
void setInstrStackFrames(bool)

\end{apient}

\apidesc{
Turn on and off stack frames in instrumentation. When on, Dyninst will create stack frames around instrumentation. A
stack frame allows Dyninst or other tools to walk a call stack through instrumentation, but introduces overhead to
instrumentation. The default is to not create stack frames.
}

\begin{apient}
bool getInstrStackFrames()

\end{apient}

\apidesc{
Return true if instrumentation will create stack frames, or false otherwise.
}

\begin{apient}
void setMergeTramp (bool)

\end{apient}

\apidesc{
Turn on or off inlined tramps. Setting this value to true will make each base trampoline have all of its
mini-trampolines inlined within it. Using inlined mini-tramps may allow instrumentation to execute faster, but
inserting and removing instrumentation may take more time. The default setting for this is true.
}

\begin{apient}
bool isMergeTramp ()

\end{apient}

\apidesc{
This returns the current status of inlined trampolines. A value of true indicates that trampolines are inlined.
}

\begin{apient}
void setSaveFPR (bool)

\end{apient}

\apidesc{
Turn on or off floating point saves. Setting this value to false means that floating point registers will never be
saved, which can lead to large performance improvements. The default value is true. Setting this flag may cause
incorrect program behavior if the instrumentation does clobber floating point registers, so it should only be used when
the user is positive this will never happen.
}

\begin{apient}
bool isSaveFPROn ()

\end{apient}

\apidesc{
This returns the current status of the floating point saves. True means we are saving floating points based on the
analysis for the given platform.
}

\begin{apient}
void setBaseTrampDeletion(bool)

\end{apient}

\apidesc{
If true, we delete the base tramp when the last corresponding minitramp is deleted. If false, we leave the base tramp
in. The default value is false.
}

\begin{apient}
bool baseTrampDeletion()

\end{apient}

\apidesc{
Return true if base trampolines are set to be deleted, or false otherwise.
}

\begin{apient}
void setLivenessAnalysis(bool)

\end{apient}

\apidesc{
If true, we perform register liveness analysis around an instPoint before inserting instrumentation, and we only save
registers that are live at that point. This can lead to faster run-time speeds, but at the expense of slower
instrumentation time. The default value is true.
bool livenessAnalysisOn()
Return true if liveness analysis is currently enabled.
}

\begin{apient}
void getBPatchVersion(int &major, int &minor, int &subminor)

\end{apient}

\apidesc{
Return Dyninst\'s version number. The major version number will be stored in major, the minor
version number in minor, and the subminor version in subminor. For example, under Dyninst 5.1.0, this function will
return 5 in major, 1 in minor, and 0 in subminor.
}

\begin{apient}
int getNotificationFD()

\end{apient}

\apidesc{
Returns a file descriptor that is suitable for inclusion in a call to select(). Dyninst will write data to this file
descriptor when it to signal a state change in the process. BPatch::pollForStatusChange should then be called so that
Dyninst can handle the state change. This is useful for applications where the user does not want to block in
BPatch::waitForStatusChange. The file descriptor will reset when the user calls BPatch::pollForStatusChange.
}

\begin{apient}
BPatch_type *createArray(const char *name, BPatch_type *ptr, unsigned int low, unsigned int hi) 

\end{apient}

\apidesc{
Create a new array type. The name of the type is name, and the type of each element is ptr. The index of the first
element of the array is low, and the last is high. The standard rules of type compatibility, described in Section
, are used with arrays created using this function.
}

\begin{apient}
BPatch_type *createEnum(const char *name, std::vector<char *> &elementNames,
                        std::vector<int> &elementIds) 

BPatch_type *createEnum(const char *name, std::vector<char *> &elementNames)

\end{apient}

\apidesc{
Create a new enumerated type. There are two variations of this function. The first one is used to create an enumerated
type where the user specifies the identifier (int) for each element. In the second form, the system specifies the
identifiers for each element. In both cases, a vector of character arrays is passed to supply the names of the
elements of the enumerated type. In the first form of the function, the number of element in the elementNames and
elementIds vectors must be the same, or the type will not be created and this function will return NULL. The standard
rules of type compatibility, described in Section , are used with enums created using this
function.
}

\begin{apient}
BPatch_type *createScalar(const char *name, int size) 

\end{apient}

\apidesc{
Create a new scalar type. The name field is used to specify the name of the type, and the size parameter is used to
specify the size in bytes of each instance of the type. No additional information about this type is supplied. The
type is compatible with other scalars with the same name and size.
}

\begin{apient}
BPatch_type *createStruct(const char *name, std::vector<char *> &fieldNames,
                          std::vector<BPatch_type *> &fieldTypes) 

\end{apient}

\apidesc{
Create a new structure type. The name of the structure is specified in the name parameter. The fieldNames and
fieldTypes vectors specify fields of the type. These two vectors must have the same number of elements or the
function will fail (and return NULL). The standard rules of type compatibility, described in Section
, are used with structures created using this function. The size of the structure is the sum of
the size of the elements in the fieldTypes vector.
}

\begin{apient}
BPatch_type *createTypedef(const char *name, BPatch_type *ptr) 

\end{apient}

\apidesc{
Create a new type called name and having the type ptr.
}

\begin{apient}
BPatch_type *createPointer(const char *name, BPatch_type *ptr)

BPatch_type *createPointer(const char *name, BPatch_type *ptr, int size) 

\end{apient}

\apidesc{
Create a new type, named name, which points to objects of type ptr. The first form creates a pointer whose size is equal
to sizeof(void*)on the target platform where the muta­tee is running. In the second form, the size of the pointer is
the value passed in the size parameter.
}

\begin{apient}
BPatch_type *createUnion(const char *name, std::vector<char *> &fieldNames,
                         std::vector<BPatch_type *> &fieldTypes) 

\end{apient}

\apidesc{
Create a new union type. The name of the union is specified in the name parameter. The fieldNames and fieldTypes
vectors specify fields of the type. These two vectors must have the same number of elements or the function will fail
(and return NULL). The size of the union is the size of the largest element in the fieldTypes vector.
}
